\section{Models and Checks}
\label{sec:setup}

\paragraph{Models, from simplest to richest.}
We organize our own models in stages. \textbf{Stage~0} holds the classical
forecasts, namely persistence ($\hat v = v_{\mathrm{pre}}$), the
training-set mean of the target (a constant), EWMA with
$\lambda{=}0.94$ \citep{riskmetrics1996}, HAR-RV \citep{corsi2009har}
and GARCH(1,1) \citep{bollerslev1986}. \textbf{Stage~1} is a LightGBM model
\citep{lightgbm2017} on past-volatility features, transcript length,
number of Q\&A turns and a \textbf{company fixed effect}, a price-and-metadata
baseline with ticker features. The company-only model of the identity checks is
separate. \textbf{Stage~2} uses the transcript through BGE-M3 \citep{bge2024}
embeddings pooled per section (1024-d), FinBERT \citep{finbert}
sentiment aggregates and surface statistics. \textbf{Stage~3} uses the audio,
resampled to 16\,kHz mono. Its features are eGeMAPSv02 functionals
\citep{egemaps2016} extracted with openSMILE 2.6.0 \citep{eyben2010opensmile},
and WavLM-Large \citep{wavlm} and emotion2vec+ \citep{emotion2vec} embeddings
over 30-second windows, averaged per call.
\textbf{Stage 4} combines modalities by
gated fusion or by stacking (a ridge over the text, audio and Stage-1 predictions,
fit on the validation set). From Stage~2 on, every model is reported with and
without a block of past-volatility features ($v_{\mathrm{pre}}$ at all horizons
plus the HAR-RV inputs). The prediction heads are a ridge regression, with its
penalty chosen on validation, and a small MLP run with five seeds. The MLPs are
unstable across seeds (standard deviation of \Rtwo{} up to 6.8), so the main grid
uses one deterministic model per stage, chosen by a rule recorded in the decision
log. MLP cells are in Appendix~\ref{app:complete}.

\paragraph{Metrics and tests.}
We report MSE on log realized volatility and, as the headline metric, \Rtwo{} relative to persistence, following \citet{gukellyxiu2020} with persistence in place of their zero forecast, defined as
$1-\sum(y-\hat y)^2/\sum(y-\hat y_{\mathrm{persist}})^2$. Because $\Delta v$ has a positive mean on the measured anchor, persistence (which predicts $\Delta v{=}0$) is a weak reference there, so every $\Delta v$ claim is also compared with the constant model. Diebold--Mariano tests \citep{dm1995} with the small-sample correction of \citet{hln1997} compare per-call squared errors against persistence and against the
previous stage. The confirmatory comparisons (each stage against Stage~1 on $\Delta v$ across horizons, Stage~4 against the best single modality, and the
published split against the controlled split for each reproduced model) are
Holm-corrected \citep{holm1979}. The central question is whether call content improves on the same
model without it, so each content model is also compared, on the same calls,
with the past-volatility ridge it extends. We report the relative MSE reduction
with a 95\% bootstrap interval that resamples companies and read it against a 2\%
margin fixed before the comparison was run. An interval above the margin is an
improvement, one below it shows no meaningful benefit, and anything else is
inconclusive. All other analyses are exploratory.

\paragraph{Identity checks.}
\label{sec:controls}
We use four. A \emph{company-only} model predicts each company's training mean
and shows how far identity alone gets. Two \emph{swaps} replace each test call's
transcript or audio, either with another call by the same company or with a
random call from the corpus. If the same-company swap leaves a model's score
unchanged, the model was reading identity, and if the random swap also leaves it
unchanged, the model was reading nothing. A linear \emph{probe} tries to recover
the ticker from the frozen call embedding. Finally, every table shows
company-disjoint results next to temporal ones. The command \texttt{ecvol audit
predictions} runs the prediction-side versions of the company-only and swap
checks on any file with one row per (call, ticker, date, split, horizon, truth,
prediction), and reports split integrity, seen versus unseen companies, the
between-company share of prediction variance and company-clustered
significance. That is a descriptive diagnostic (Appendix Table~\ref{tab:toolkit}). A permutation of existing predictions moves
the price inputs along with the content, so it is not a substitute for the
input-level swaps, which need the model itself. The price-only ridge of
Table~\ref{tab:grid} scores the same under a real transcript and under either
swap, as it must.

\paragraph{Reproductions and post-cutoff tests.}
Of seventeen papers in this line, one releases code that runs as released
(Same-Company-Same-Signal), seven release code that needs repairs or inputs that
were never shipped, and nine release nothing usable (Table~\ref{tab:code-audit}).
We rerun the five with usable code (HTML, KeFVP, Same-Company-Same-Signal,
\citealp{sawhney2020mm}, DialogueGAT) \emph{on their own data}, using the authors'
code wherever it can be made to start and a line-by-line port where it cannot,
with every patch in a released ledger. We compare with the published number and
then rerun the same model under a gapped and a \tickerdisjoint{} split of the same
data. Earnings25 allows two more tests, the frozen Stage-3 audio models rerun by
bitrate, and the whole FinCall-trained pipeline scored, with no retraining, on
calls that postdate every model's training cutoff.
